\chapter{\textlatin{Scripts} Επίδειξης}
\label{app:demo-scripts}

Το παρόν παράρτημα περιέχει τα \textlatin{scripts} που χρησιμοποιήθηκαν
για την επίδειξη του κύκλου εκχώρησης πρόσβασης. Τα \textlatin{scripts}
εκτελούνται σε \textlatin{PowerShell} και αξιοποιούν τα εργαλεία
\textlatin{curl} και \textlatin{Docker}.

% ------------------------------------------------------------------
\section{Επίδειξη Παρόχου (\textlatin{demo.ps1})}
\label{app:demo-provider}
% ------------------------------------------------------------------

Το \textlatin{script} \textlatin{demo.ps1} εκτελεί δύο σενάρια: πρόσβαση
χωρίς \textlatin{token} (αναμένεται \textlatin{401}) και πρόσβαση
με \textlatin{token} παρόχου (αναμένεται \textlatin{200}) για τρεις
τύπους οντοτήτων.

\selectlanguage{english}
\begin{lstlisting}[language={},]
# demo.ps1
# =========================================================
# Access control demonstration via Kong/PEP with iSHARE JWT
#
# Scenario 1: No token        -> 401 Unauthorized
# Scenario 2: Provider token  -> 200 OK + data
# =========================================================

$ISHARE_DIR = "C:\Users\SAKak\Documents\GitHub\Data-Spaces\Kong\iShare"
$KONG_BASE  = "http://localhost:8000/ngsi-ld/v1/entities"

$ELECTRIC_TYPE = "https://plegma.example.org/vocab%23HouseholdElectricMeasurement"
$ENV_TYPE      = "https://plegma.example.org/vocab%23EnvironmentalMeasurement"
$BUILDING_TYPE = "https://smartdatamodels.org/dataModel.Building/Building"

Write-Host "============================================================"
Write-Host "  Plegma Data Space - Provider Access Control Demo"
Write-Host "============================================================"

# 1. No token -> expect 401
Write-Host "`n=== 1. No token (expected: 401) ==="
curl.exe -s "$KONG_BASE`?type=$ELECTRIC_TYPE&limit=1"

# 2. Provider token -> generate fresh JWT and send request
Write-Host "`n=== 2. Provider token (expected: 200) ==="
$PROVIDER_TOKEN = docker run --rm `
    -v "${ISHARE_DIR}:/iShare" `
    python:3.9-slim sh -c `
    "pip install PyJWT cryptography -q 2>/dev/null && python3 /iShare/gen_token.py"

Write-Host "HouseholdElectricMeasurement:"
curl.exe -s -H "Authorization: Bearer $PROVIDER_TOKEN" `
    "$KONG_BASE`?type=$ELECTRIC_TYPE&limit=1" | python -m json.tool

Write-Host "`nEnvironmentalMeasurement:"
curl.exe -s -H "Authorization: Bearer $PROVIDER_TOKEN" `
    "$KONG_BASE`?type=$ENV_TYPE&limit=1" | python -m json.tool

Write-Host "`nBuilding:"
curl.exe -s -H "Authorization: Bearer $PROVIDER_TOKEN" `
    "$KONG_BASE`?type=$BUILDING_TYPE&limit=1" | python -m json.tool
\end{lstlisting}
\selectlanguage{greek}

% ------------------------------------------------------------------
\section{Επίδειξη Καταναλωτή (\textlatin{demo\_consumer.ps1})}
\label{app:demo-consumer}
% ------------------------------------------------------------------

Το \textlatin{script} \textlatin{demo\_consumer.ps1} εκτελεί τρία σενάρια:
πρόσβαση χωρίς \textlatin{token}, πρόσβαση με \textlatin{token} παρόχου,
και πρόσβαση με \textlatin{token} καταναλωτή μέσω \textlatin{delegation}.
Το τρίτο σενάριο περιλαμβάνει και τα τρία βήματα της ροής
\textlatin{M2M} του καταναλωτή: δημιουργία \textlatin{iSHARE JWT},
ανταλλαγή με \textlatin{access token} μέσω \textlatin{Keyrock} και
πρόσβαση μέσω \textlatin{Kong}.

\selectlanguage{english}
\begin{lstlisting}[language={},]
# demo_consumer.ps1
# =========================================================
# Access delegation cycle demonstration (iSHARE delegation)
#
# Scenario 1: No token              -> 401 Unauthorized
# Scenario 2: Provider token (M2M)  -> 200 OK
# Scenario 3: Consumer token (M2M)  -> 200 OK (via delegation)
#
# Consumer Flow:
#   1. Consumer creates iSHARE JWT with consumer.key
#   2. Consumer sends JWT to Keyrock -> gets access token
#      (Keyrock verifies consumer via Satellite + finds delegation policy)
#   3. Consumer sends access token to Kong -> gets data
# =========================================================

$ISHARE_DIR    = "C:\Users\SAKak\Documents\GitHub\Data-Spaces\Kong\iShare"
$KONG_BASE     = "http://localhost:8000/ngsi-ld/v1/entities"
$KEYROCK_URL   = "http://localhost:3007/oauth2/token"
$CONSUMER_EORI = "EU.EORI.NLPLEGMACONSUMER"
$ELECTRIC_TYPE = "https://plegma.example.org/vocab%23HouseholdElectricMeasurement"

Write-Host "============================================================"
Write-Host "  Plegma Data Space - Consumer Access Control Demo"
Write-Host "============================================================"

# --- Scenario 1: No token ---
Write-Host "`n=== 1. No token (expected: 401 Unauthorized) ==="
curl.exe -s "$KONG_BASE`?type=$ELECTRIC_TYPE&limit=1"

# --- Scenario 2: Provider token (M2M) ---
Write-Host "`n=== 2. Provider iSHARE JWT (expected: 200 OK) ==="
$PROVIDER_TOKEN = docker run --rm `
    -v "${ISHARE_DIR}:/iShare" `
    python:3.9-slim sh -c `
    "pip install PyJWT cryptography -q 2>/dev/null && python3 /iShare/gen_token.py"

curl.exe -s -H "Authorization: Bearer $PROVIDER_TOKEN" `
    "$KONG_BASE`?type=$ELECTRIC_TYPE&limit=1" | python -m json.tool

# --- Scenario 3: Consumer token (M2M with delegation) ---
Write-Host "`n=== 3. Consumer token via Keyrock delegation (expected: 200 OK) ==="

# Step 1: Create consumer iSHARE JWT
Write-Host "  Step 1/3: Creating iSHARE JWT for Consumer..."
$CONSUMER_JWT = docker run --rm `
    -v "${ISHARE_DIR}:/iShare" `
    python:3.9-slim sh -c `
    "pip install PyJWT cryptography -q 2>/dev/null && python3 /iShare/gen_consumer_token.py"

# Step 2: Exchange JWT for access token from Keyrock
Write-Host "  Step 2/3: Getting access token from Keyrock..."
$tokenResponse = curl.exe -s -X POST $KEYROCK_URL `
    -H "Content-Type: application/x-www-form-urlencoded" `
    -d "grant_type=client_credentials&scope=iSHARE&client_assertion_type=urn:ietf:params:oauth:client-assertion-type:jwt-bearer&client_assertion=$CONSUMER_JWT&client_id=$CONSUMER_EORI"

$ACCESS_TOKEN = ($tokenResponse | ConvertFrom-Json).access_token

if (-not $ACCESS_TOKEN) {
    Write-Host "  ERROR: No access token received from Keyrock!"
    Write-Host "  Response: $tokenResponse"
    exit 1
}
Write-Host "  Access token received (length: $($ACCESS_TOKEN.Length) chars)"

# Step 3: Use access token to access data via Kong
Write-Host "  Step 3/3: Sending request to Kong with access token..."
curl.exe -s -H "Authorization: Bearer $ACCESS_TOKEN" `
    "$KONG_BASE`?type=$ELECTRIC_TYPE&limit=1" | python -m json.tool

Write-Host "`n============================================================"
Write-Host "  Results:"
Write-Host "  1. No token       -> 401 Unauthorized"
Write-Host "  2. Provider token -> 200 OK"
Write-Host "  3. Consumer token -> 200 OK (via delegation)"
Write-Host "============================================================"
\end{lstlisting}
\selectlanguage{greek}
